% !TEX TS-program = xelatex
% !TEX encoding = UTF-8 Unicode
% !BIB TS-program = bibtex

% Integrable structure as a learnable inductive bias -- English version (draft).
% Build: ./build.sh  or  latexmk -xelatex main.tex
% Parallel translation of ../latex_cn/main.tex.

\documentclass[11pt]{article}

% --- math / fonts --------------------------------------------------------
\usepackage{amsmath, amssymb, amsfonts, amsthm}
\usepackage{mathrsfs}

% --- citations / bibliography --------------------------------------------
\usepackage[numbers]{natbib}

% --- tables / figures ----------------------------------------------------
\usepackage{array}
\usepackage{booktabs}
\usepackage{tabularx}
\usepackage{graphicx}
\usepackage{float}
\usepackage{subcaption}

% --- layout --------------------------------------------------------------
\usepackage[margin=0.9in]{geometry}
\usepackage{verbatim}

% --- hyperlinks (load last) ----------------------------------------------
\usepackage[colorlinks=true, citecolor=blue, linkcolor=blue, urlcolor=blue,
            hypertexnames=false]{hyperref}

% --- theorem environments ------------------------------------------------
\newtheorem{Theorem}{Theorem}[section]
\newtheorem{Definition}[Theorem]{Definition}
\newtheorem{Proposition}[Theorem]{Proposition}
\newtheorem{Lemma}[Theorem]{Lemma}
\newtheorem{Remark}[Theorem]{Remark}
% Empirical results (numerical findings, not mathematical theorems) use an
% upright style, so that experiments are not dressed up as theorems.
\theoremstyle{definition}
\newtheorem{Result}{Result}
\theoremstyle{plain}

% --- metadata ------------------------------------------------------------
% TODO (final): title still open. Current working title.
\title{Integrable Structure as a Learnable Inductive Bias:\\
       An End-to-End Case Study for a Physics Model Library}
\author{Dong Zhang\thanks{Email: dongzhanghz@gmail.com}}
\date{\today}

\begin{document}
\maketitle

% --- custom macros -------------------------------------------------------
\newcommand{\amp}{\mathrm{amp}}

% ============ writing guardrails (read before editing; delete before final) ====
%  Guardrail 1 (boundary / related work, or it gets rejected in one line):
%    - "recurrent / state-tracking models extrapolate in length, attention does not"
%      is an established result (Liu 2023; Merrill 2023/2024). We do NOT sell
%      "beating Transformers".
%    - Physics inductive biases have precedent (HNN Greydanus 2019; LNN Cranmer 2020;
%      invertible nets NICE/RealNVP). State the delta: those mostly treat physics as a
%      target or as a guarantee; we measure whether integrable structure lets a network
%      *learn* interaction content that neither a generic model nor a bolt-on can, and
%      keep it honest with the leak / bolt-on audits.
%    - The Box-Ball System is an established object in mathematical physics
%      (Takahashi-Satsuma 1990); no one has used it as a learnable sequence model.
%  Guardrail 2 (scope, do not cross):
%    - Claim only the "discrete integrable" case. Toda (continuous) is an honest
%      negative, written up in the boundary section, used to draw the line.
%    - The point is exactness / interaction content, not conservation -- conservation
%      is bolt-on-able (shown by the bolt-on control).
%    - Do not say "a full benchmark is built": this is a validated Tier 1 (one system, K=3).
%  Guardrail 3 (honest labels): hard numerical (multi-seed) / illustrative / outlook.
%    The two audits (leak, bolt-on) are the honesty method; every main table carries them.
% ==================================================================

\begin{abstract}
\noindent
% TODO: compress once more before final. Current version states scope honestly.
Neural network architectures are usually found by empirical search: a design survives
local modification, cross-task transfer, and benchmark trial and error, rather than
being selected from a systematic library of candidate structures. This paper treats
physics as one such library of candidate models and works through a single entry,
integrable systems, from end to end. We start from a physical structure, turn it into
a trainable network, design a task that tests its structural advantage, and use control
audits to rule out hard-coding and bolt-on shortcuts. Concretely, we report a positive
result on a \textbf{discrete} integrable system. On a \emph{finite-carrier Box-Ball
System}, a model built around the integrable carrier structure (a conservation-respecting
carrier sweep with a learnable emit gate) learns to preserve the \textbf{soliton-amplitude
content} exactly across long compositions of many collisions, an invariant richer than
the ball count. \emph{Under the models, training budget, and evaluation protocol we test},
generic sequence models (GRU, LSTM, Transformer) do not, and bolting the conserved
quantity onto a generic model does not either. Two audits keep the claim honest. A
\textbf{leak audit} (replacing the learnable gate with a fixed rule) shows the advantage is
\emph{learned}, not hard-coded: the fixed rule reaches only $40\%$ amplitude content, the
learned gate $100\%$. A \textbf{bolt-on control} shows the advantage \emph{cannot} be
bolted on: pinning the ball count to $100\%$ recovers only $0.7\%$ amplitude content, so
conservation transfers while the interaction content does not. All results hold across
$5$ random seeds. We also mark the boundary: the effect does not carry over to a smooth
continuous system (the Toda lattice is a negative result), so we claim only the discrete
integrable case. We do not present the Box-Ball System as a general architecture; the
contribution is a demonstration that the ``physics model library'' route can produce a
measurable, cheat-resistant machine-learning result.
\end{abstract}

% \tableofcontents

% ========================================================================
\section{Introduction}\label{sec:intro}
% ========================================================================

The evolution of deep learning models looks more like accumulated empirical engineering
than a selection process driven by first principles. This does not mean the existing
architectures were built without rational motivation. In vision, the local connectivity
and weight sharing of CNNs have a clear motivation~\citep{lecun1998gradient}, yet the
breakout of AlexNet depended on ImageNet, GPUs, and specific engineering
tricks~\citep{krizhevsky2012imagenet}. In language, RNNs~\citep{elman1990finding},
LSTMs~\citep{hochreiter1997long}, Transformers~\citep{vaswani2017attention},
BERT~\citep{devlin2019bert}, and the GPT series~\citep{brown2020language} answered, in turn,
sequential recurrence, long-range dependence, parallelizable attention, pretraining, and
scale. In generative modeling, GANs~\citep{goodfellow2014generative},
VAEs~\citep{kingma2014autoencoding}, and diffusion
models~\citep{sohldickstein2015deep, ho2020denoising} came from adversarial training,
variational inference, and stepwise noising and denoising. Each model has a local
motivation, yet these local reasons do not add up to a systematic methodology for building
models. A new architecture is usually proposed under a specific bottleneck, a modeling
intuition, or benchmark pressure, and then filtered by data scale, compute, training
stability, benchmark numbers, and community reuse~\citep{lipton2018troubling}. The gap is in
model discovery, not model design. Design has clear reasons; discovery has no systematic
space of candidates to draw from. Where a candidate structure comes from, which tasks it
should be tried on first, and how its failure boundaries accumulate all tend to rely on
individual experience and community trial and error~\citep{elsken2019neural}. In practice,
building a model still means editing an existing architecture, porting a familiar structure
to a new setting, or assembling a design by intuition and handing it to a
benchmark~\citep{he2021automl,radosavovic2020designing}. What is missing is a library of
candidate structures that lets us select and try models systematically and turn the
modeling process into engineering; more one-off inspiration would not fill that gap.

This is the gap the research program behind this paper aims at. Machine learning today does
not lack model families. Tree-based models stay strong on many tabular
tasks~\citep{grinsztajn2022tree}, CNNs still serve a large range of vision and spatial
problems, and Transformers and diffusion models carry the mainstream of language and
generation. In science, domain models from protein structure prediction (AlphaFold,
RoseTTAFold) to weather forecasting (GraphCast) show how much a structural prior can decide a
specific scientific task~\citep{jumper2021highly,baek2021accurate,lam2023learning}. This
coexistence of many families points to a different question: how to generate, select, and
accumulate candidate structures systematically when a new task appears, rather than how to
propose one more model. The current ecosystem mostly stores results that have already been
invented and validated, instead of a mechanism for generating and selecting candidates for
future tasks. The match between a model and a task is likewise scattered across papers,
leaderboards, community lore, and individual
intuition~\citep{lipton2018troubling,elsken2019neural,he2021automl,radosavovic2020designing}.
What we ultimately want to build is an infrastructure of candidate spaces for model
discovery, not a list of model names, a model zoo, or a leaderboard. Such a library has to
record three things: where a candidate structure comes from, what task features it suits,
and where it fails. Once this information can be searched, compared, and accumulated, model
discovery grows into a reusable, comparable, and extensible engineering process.

Our proposal is that \textbf{physics itself is such a library of candidate models, open to
systematic mining}. Physics has inspired neural networks before. The clearest early line
runs through statistical physics: the energy function, state flips, and low-energy structure
of the Ising model and spin glasses directly inspired the Hopfield network, and the Boltzmann
distribution and thermal-equilibrium picture entered the Boltzmann machine and the restricted
Boltzmann machine~\citep{hopfield1982neural,hinton1985learning,smolensky1986information}. The
physics here is not decorative analogy; it entered the network itself as energy functions,
probability distributions, and sampling mechanisms. That line also exposed a typical limit of
physics-inspired models: energy models are often hard to train and expensive to sample, and
they do not scale naturally to large data and modern deep architectures. Physical inspiration
does not automatically yield a working architecture; it has to be re-engineered and then prove
its advantage on a concrete task.

Diffusion models are a more recent example. The idea first drew on the ``add noise step by
step, then reverse'' picture from non-equilibrium thermodynamics, writing generative modeling
as a process from data to noise and back~\citep{sohldickstein2015deep}. After later
simplification and engineering, that physical picture became the standard training and
sampling procedure of today's deep generative models~\citep{ho2020denoising}. The physics
label can recede into the background as the method matures, while the process structure it
supplied stays in place. Beyond this, symmetry and equivariance inspired CNNs, GNNs, and
geometric deep learning, and Hamiltonian or oscillator structure has been used for dynamical
systems and long-range sequence
modeling~\citep{bronstein2021geometric,greydanus2019hamiltonian,rusch2021coupled,rusch2024oscillatory}.
Physics has entered neural networks many times, as energy, distributions, dynamical
processes, symmetry, and structural constraints.

These successes did not follow one clean path. The Ising-and-Boltzmann-machine line
contributed core concepts to deep learning, such as energy models and stochastic units, yet
as a directly trainable, deployable architecture it did not grow into today's deep learning
backbone. Diffusion models tell a different story: the physical idea was proposed by
Sohl-Dickstein et al.\ as early as 2015~\citep{sohldickstein2015deep}, stayed quiet for
several years, and was reactivated around 2020 by the engineering of DDPM and related
work~\citep{ho2020denoising} to become mainstream in generative modeling. This history shows
that the value of a physics model to AI is not an abstract ``useful'' or ``useless''; it
depends on whether the model can be rewritten into a trainable structure, whether it matches
some class of tasks, and where its failure boundary lies.

So the question we care about is not to add a few more examples of physics inspiring AI, but
whether we can make one thing systematic: take a physics model, especially the many mature
many-body models physics offers, as a source of candidate structures; rewrite it into a
trainable network; let it learn directly from task data, as a physics model turned into a
learning engine; then test it on tasks, and record what it suits and where it fails. Put
differently, the ``trial-and-error search for a model'' that once relied on scattered
individual intuition becomes a process of sourcing candidates from physics and testing them
systematically. Physics happens to supply the raw material for this. The Ising model, the
Boltzmann distribution, and non-equilibrium diffusion mentioned above use only a small part
of the physics toolbox; beyond them, physics has accumulated many mature models and
mathematical skeletons, including integrable systems, open systems, many-body interactions,
conservation laws, symplectic structure, reaction-diffusion, and multiscale theory. Many-body
models in particular form a large and organized supply of candidate structures (Appendix~\ref{app:catalog}
collects these physics families with their corresponding AI capabilities and maturity in one
table). The real question follows: can these models form a library of candidate models for AI?
Can we answer, in a standardized way and one by one, what state space each corresponds to, how
it becomes a trainable network, what task it suits, where its advantage over existing models
lies, and under what conditions it fails?

This paper does not try to build that whole program. It makes a proof of concept first. In this
proof of concept we pick one physics model that has not been used in deep learning before, turn
it into a trainable network, put it on a benchmark, and see under what conditions it is actually
useful, and even more useful than the currently dominant deep learning models, so that the
questions above get answered once on one concrete model. Along the way we start from the
physical structure, design a task that genuinely tests its structural advantage, use control
audits to rule out hard-coding and bolt-on shortcuts, and mark its capability boundary. The
first entry we mine is the \textbf{integrable / reversible system}. Its relevance is not
``using AI to solve physics problems'', but that it offers an unusual prototype of sequence
interaction: the system can undergo strong nonlinear interaction while the interaction content
is preserved \emph{exactly} (solitons collide, pass through, and shift phase, and their
amplitude content stays intact). A generic neural sequence model, by contrast, can learn
approximate dynamics, but under long composition and length extrapolation a small per-step
error accumulates, and the conserved quantities, reversibility, and richer structural content
drift. So the real test is not whether a hard-coded integrable rule, treated as a model, can
extrapolate (that is circular), but rather: \textbf{can a neural model that carries an
integrable inductive bias, yet still has to learn the key rule from data, learn the interaction
content that neither a generic model nor a bolt-on constraint can reach?}

Our proof of concept lands on the \textbf{finite-carrier Box-Ball System}. A plain Box-Ball
System leaks easily to hard-coding: a fixed rule that ignores the carrier already solves the
task. Once the carrier capacity is limited to $K$, the rule becomes nontrivial: when the
carrier is full, a ball has to pass through, and the correct drop decision has to use the
carrier state. We also go beyond ball-count conservation and evaluate the \textbf{soliton-amplitude
content}, the multiset of soliton amplitudes. It runs deeper than the ball count, because the
same ball count can decompose into different amplitudes (for example $\{5,3,1\}$ and $\{4,4,1\}$
are both $9$ balls), so bolting ball-count conservation onto any model cannot recover it. On the
task with $K=3$ and composition $T=8$, the model that builds in this integrable carrier structure
reaches $100\pm0\%$ per-cell accuracy, ball-count conservation, and amplitude-content preservation
across $5$ random seeds. The two audits give the key evidence: replacing the learnable emit gate
with a fixed rule drops amplitude content to $40\pm3\%$, so the $60$-point advantage comes from
learning rather than hard-coding; projecting a generic model's output onto the correct ball count
lifts ball-count conservation to $100\%$ but leaves amplitude content at $0.7\pm0.8\%$, so the
advantage is not something a scalar conservation constraint can reproduce (Table~\ref{tab:tier1}).
We mark the boundary at the same time: this is a positive result on a \textbf{discrete} integrable
system, not a full benchmark family; on the smooth continuous Toda lattice the same discriminating
effect does not appear (\S\ref{sec:boundary}).

The boundary with existing work needs stating. That ``recurrent / state-tracking models with a
left-to-right sweep extrapolate in length while attention does not'' is an established
result~\citep{liu2023transformers, merrill2023parallelism, merrill2024illusion}, and we do
\emph{not} claim ``beating Transformers'' as a contribution; the Transformer is a structure-free
reference here. Using physical structure as an inductive bias also has precedent (Hamiltonian
networks~\citep{greydanus2019hamiltonian}, Lagrangian networks~\citep{cranmer2020lagrangian},
invertible coupling networks~\citep{dinh2014nice, dinh2016density}), but those mostly treat
physics as a \emph{target} or as a \emph{guarantee}. Our difference is to measure whether
integrable structure lets a network \emph{learn} interaction content that neither a generic
model nor a bolt-on constraint can reach, with the leak and bolt-on audits nailing down ``learned
vs.\ hard-coded'' and ``structure vs.\ conservation''. The Box-Ball System is itself an
established integrable object in mathematical physics~\citep{takahashi1990soliton}, but as far as
we know no one has used it as a learnable sequence model and compared it against mainstream models
under length / composition extrapolation.

To sum up, this paper makes four contributions. (i) It turns the ``physics model library'' route
into an end-to-end, reproducible, cheat-resistant machine-learning case study rather than a
manifesto. (ii) On the finite-carrier Box-Ball System it constructs a learnable model with
integrable carrier structure and gives evidence that it preserves soliton-amplitude content
exactly under long composition (\S\ref{sec:main}). (iii) It introduces the leak audit and the
bolt-on control as two methodological guardrails, separating ``learned vs.\ hard-coded'' and
``integrable content vs.\ scalar conservation'' (\S\ref{sec:setup:audits}). (iv) It states the
boundary honestly: the result is limited to the discrete integrable case for now, and the
continuous Toda negative result shows the mechanism does not transfer automatically
(\S\ref{sec:boundary}).

% ========================================================================
\section{Background and setup}\label{sec:setup}
% ========================================================================

\subsection{The Box-Ball System and its finite-carrier variant}\label{sec:setup:bbs}
The Box-Ball System (BBS) is a discrete integrable cellular automaton on a one-dimensional
$0/1$ lattice: a \emph{carrier} sweeps from left to right, picking up balls it meets and
dropping them into empty cells. A run of consecutive $1$s is a \textbf{soliton}; its amplitude
equals the run length, it moves right at a speed proportional to its amplitude, and when a large
soliton catches a small one it \emph{passes through} with each amplitude unchanged. The
\textbf{finite-carrier Box-Ball System} caps the carrier capacity at $K$: when the carrier is
full (holding $K$ balls), a ball it meets \emph{passes through} instead of being picked up. This
change makes the single-step rule \emph{nontrivial}: the correct drop decision has to use the
current carrier count, so a fixed rule that ignores the carrier cannot solve it (this is the basis
of our leak audit, \S\ref{sec:setup:audits}).
\textbf{TODO: add a single-step schematic or pseudocode for capacity $K$.}

\subsection{Interaction content: an invariant richer than the ball count}\label{sec:setup:content}
The total ball count is the shallowest conserved quantity, a scalar; it \emph{can} be bolted onto
any model (\S\ref{sec:main}). The signature of integrability is much richer: the \textbf{multiset
of soliton amplitudes} $\amp(x)=\{a_1\ge a_2\ge\cdots\}$ stays \emph{exactly invariant} under the
dynamics, and it \textbf{cannot} be recovered from the ball count, since two states can share a
ball count while differing in amplitude multiset (for example $\{5,3,1\}$ and $\{4,4,1\}$ are both
$9$ balls). We use this multiset as the measure of whether the integrable interaction content has
been preserved. It is read off by evolving the state until the solitons fully separate and reading
the run lengths; we checked numerically that it is conserved under the finite-carrier dynamics and
that the system is reversible ($K=2/4/6$).
% Maturity: hard numerical (self-check passes).

\subsection{Two audits: the honesty method for this paper}\label{sec:setup:audits}
\begin{Definition}[leak audit / carrier-blind]\label{def:leak}
Replace the \emph{learnable} emit gate in the structural model with a \emph{fixed} rule
(\texttt{emit}$=1-$cell, untrained), keeping everything else. If the fixed version scores about the
same as the learned one, the ``learning'' is empty and the advantage comes from hard-coding; if the
fixed version is clearly worse, the advantage is \emph{genuinely learned}.
\end{Definition}

\begin{Definition}[bolt-on control]\label{def:bolton}
Take a generic model and, at inference, \emph{project} its output onto the constraint ``same ball
count as the input'' (keep the $N$ highest-probability cells, with $N=$ input ball count). This
\emph{forces} the scalar conserved quantity onto the model. If it still fails to recover the
interaction content, then that content \emph{cannot} be bolted on and is specific to the structure.
\end{Definition}

% ========================================================================
\section{Main result: integrable extrapolation on the finite-carrier BBS (Tier 1)}\label{sec:main}
% ========================================================================

\paragraph{Task.}
Learn \emph{one} finite-carrier BBS step, then \textbf{compose} it on multi-soliton states through
many collisions (train on few, short solitons; test on many, long ones). We evaluate four quantities:
per-cell accuracy, ball-count conservation, and the \emph{exact-match rate} and \emph{IoU} of the
soliton-amplitude content. Entrants: the structural model (a conservation-respecting carrier sweep
with a learnable emit gate, \S\ref{sec:setup:bbs}); generic models GRU / LSTM / Transformer; the
bolt-on control (GRU $+$ ball-count projection, Definition~\ref{def:bolton}); and the leak audit
(carrier-blind, Definition~\ref{def:leak}).

\begin{table}[h]
\centering
\caption{Integrable-extrapolation leaderboard on the finite-carrier Box-Ball System ($K=3$,
composition $T=8$ steps, $5$ random seeds, mean $\pm$ std). The structural model is $100\pm0$ on all
four metrics; the leak audit reaches only $40\pm3$ amplitude content (a $60$-point learning margin,
no leak); the bolt-on control pins the ball count at $100\%$ but recovers only $0.7\%$ amplitude
content (conservation transfers, content does not).}
\label{tab:tier1}
\begin{tabular}{lcccc}
\toprule
\textbf{Entrant} & \textbf{Acc.\%} & \textbf{Ball-cons.\%} & \textbf{Amp.-exact\%} & \textbf{Amp.-IoU\%} \\
\midrule
\textbf{Structural carrier model} & $\mathbf{100.0\pm0.0}$ & $\mathbf{100.0\pm0.0}$ & $\mathbf{100.0\pm0.0}$ & $\mathbf{100.0\pm0.0}$ \\
GRU (generic) & $91.7\pm1.2$ & $0.0\pm0.0$ & $0.0\pm0.0$ & $5.9\pm11.2$ \\
LSTM (generic) & $92.7\pm1.3$ & $1.7\pm3.3$ & $0.0\pm0.0$ & $13.5\pm11.6$ \\
Transformer (generic) & $71.6\pm10.9$ & $0.0\pm0.0$ & $0.0\pm0.0$ & $20.3\pm3.6$ \\
Bolt-on (GRU $+$ pinned count) & $89.2\pm2.5$ & $100.0\pm0.0$ & $0.7\pm0.8$ & $21.8\pm5.3$ \\
Leak audit (carrier-blind) & $88.9\pm0.5$ & $100.0\pm0.0$ & $40.0\pm3.2$ & $57.9\pm2.3$ \\
\bottomrule
\end{tabular}
\end{table}

\begin{Result}[Learned, not hard-coded]\label{res:genuine}
% Maturity: hard numerical (5 seeds). This is an empirical finding, not a theorem.
The structural model's advantage is \emph{learned}: the leak audit (fixed emit gate) reaches only
$40.0\pm3.2\%$ amplitude content, while the learned gate reaches $100\%$, a $60$-point learning
margin. This is not the leak of the plain BBS (where the fixed rule also scores $100\%$); here the
fixed rule is clearly worse, which shows the emit gate really did learn to use the carrier state.
\end{Result}

\begin{Result}[Deep, cannot be bolted on]\label{res:depth}
% Maturity: hard numerical (5 seeds).
Bolting the ball count onto a generic model lifts ball-count conservation to $100\%$, but recovers
only $0.7\pm0.8\%$ amplitude content. The scalar conserved quantity can be bolted on; the soliton
amplitude multiset cannot. \emph{Under the models, training budget, and evaluation protocol we test},
only the model with integrable carrier structure keeps it, and the generic models (without bolt-on)
collapse on both (amplitude content $0\%$).
\end{Result}

% ========================================================================
\section{Mechanism}\label{sec:mechanism}
% ========================================================================
\textbf{TODO: main text.} Three points:
(i) structural conservation (carrier bookkeeping) constrains the single-step rule to be
conservative, so it can be learned \emph{exactly}, and composing it many times does not drift;
(ii) a generic model only learns an approximate single step, and discrete composition amplifies the
small error, losing content first and ball count later;
(iii) a bolt-on only pins a scalar constraint and cannot pin the richer amplitude multiset (the
amplitude pairing after collisions).

% ========================================================================
\section{Supporting experiments}\label{sec:support}
% ========================================================================

\subsection{Box-Ball (finite-carrier): exactness beats free-form, conservation bolts on}\label{sec:support:bbs}
% Maturity: hard numerical (multi-seed). See experiments/box_ball_learned_vs_transformer.
In a length-extrapolation setting (train $L=32$, test up to $L=256$, $5$ seeds), the structural
carrier model extrapolates exactly at $100\%$, while generic sweep models (GRU / LSTM / SSM) drift
with length; the control that bolts conservation onto a free-form model can pin the conservation rate
back to $100\%$, yet \emph{cannot} recover the accuracy (Table~\ref{tab:boxball}, a gap of $9$--$21$
points). This confirms the core distinction of the main result: \textbf{conservation bolts on,
exactness does not}.

\begin{table}[h]
\centering
\caption{Finite-carrier Box-Ball ($K=6$) length extrapolation, per-cell accuracy / ball-count
conservation at $L=256$ (OOD), $5$ seeds. The structural model is $100/100$; free-form models drift
with length; the bolt-on control pins conservation back to $100$ but stays at $91$ accuracy (a $9$-point
gap). The Transformer has no sweep prior and is a reference only. See
\texttt{experiments/box\_ball\_learned\_vs\_transformer}.}
\label{tab:boxball}
\begin{tabular}{lcc}
\toprule
\textbf{Entrant} & \textbf{Acc.\%} & \textbf{Ball-cons.\%} \\
\midrule
\textbf{Structural carrier model} & $\mathbf{100}$ & $\mathbf{100}$ \\
LSTM (free-form) & $91$ & $6$ \\
GRU (free-form) & $79$ & $0$ \\
Transformer (reference) & $64$ & $0$ \\
Bolt-on (LSTM $+$ pinned conservation) & $91$ & $100$ \\
\bottomrule
\end{tabular}
\end{table}

\subsection{Margolus block automaton: the effect survives a different reversible backbone}\label{sec:support:margolus}
% Maturity: hard numerical (multi-seed + bolt-on control). See experiments/margolus_block_ca.
The same ``structural conservation $+$ composition'' recipe carries over to a structurally different
reversible system (a Margolus block-partitioning CA, with no long-range carrier; the step count grows
with length, $T\propto L$). Even when a free-form model learns the single step to $\sim99\%$, long
composition collapses its conservation to $\sim0$ and its accuracy to chance; bolting conservation on
pins the conservation rate back but still cannot recover the accuracy; the structural model stays
$100/100$ throughout (Table~\ref{tab:margolus}). So the main result is not specific to the
Box-Ball / carrier setting.

\begin{table}[h]
\centering
\caption{Margolus block CA ($5$ seeds): single-step accuracy, and accuracy / conservation composed to
$L=384$. The free-form models nearly saturate the single step, then collapse on conservation and drop
to chance accuracy under long composition; the bolt-on control pins conservation back to $100$ without
recovering accuracy; the structural block-CA is $100/100$ throughout. See
\texttt{experiments/margolus\_block\_ca}.}
\label{tab:margolus}
\begin{tabular}{lccc}
\toprule
\textbf{Entrant} & \textbf{Single-step acc.\%} & \textbf{Composed acc.\% ($L{=}384$)} & \textbf{Composed cons.\% ($L{=}384$)} \\
\midrule
\textbf{Structural block-CA} & $\mathbf{100}$ & $\mathbf{100}$ & $\mathbf{100}$ \\
bi-LSTM (free-form) & $98.7$ & $53$ & $1$ \\
Transformer (free-form) & $99.7$ & $50$ & $0$ \\
Bolt-on (bi-LSTM $+$ pinned cons.) & $98.7$ & $56$ & $100$ \\
\bottomrule
\end{tabular}
\end{table}

% ========================================================================
\section{Boundary: discrete $\to$ continuous does not transfer (Toda negative)}\label{sec:boundary}
% ========================================================================
% Maturity: numerical (negative result).
The mechanism behind the discriminating effect in this paper \emph{depends on discreteness}. On the
Toda lattice (genuinely integrable, but \textbf{smooth and continuous}), a free-form composing stepper
\emph{already} solves the time-extrapolation task well enough (state error $\sim0.003$), so the
structure has no gap to win: the small step of a smooth flow is captured exactly by the free-form
model, composition does not accumulate, and there is no discriminating mechanism like the discrete
BBS's ``off by a little each step, amplified to disaster as $T\propto L$''. We record Toda honestly as
a negative result, and on that basis \textbf{restrict the claim of this paper to the discrete integrable
case}. We also tried two ways to \emph{learn} the integrable structure (a learnable action-angle model
and a Neural-Lax model); neither beat the bolt-on baseline in the continuous setting, and both are
recorded as negative results (Appendix~\ref{app:negatives}).

% ========================================================================
\section{Related work}\label{sec:related}
% ========================================================================
\textbf{TODO: three layers, with the delta stated explicitly (see guardrail 1).}

% ========================================================================
\section{Discussion and outlook}\label{sec:outlook}
% ========================================================================
This paper gives a positive result on a discrete integrable system, cheat-resistant (two audits) and
confirmed across multiple seeds: integrable structure lets a network learn interaction content that
neither a generic model nor a bolt-on can reach. This is the first brick in grounding ``an integrable
physics model as a learnable inductive bias'', not a complete benchmark family.
\textbf{Outlook (deferred, not claimed here):} extend the validated Tier 1 into a family across several
discrete integrable / reversible systems and capacities (including a colored Box-Ball System with
labels, which makes the invariant richer); and the deeper question of making ``continuous integrable
(KdV / Lax)'' learnable, which stays open (the Toda negative gives it an honest lower bound).

\appendix
% ========================================================================
\section{A catalog of physics candidate model families (the mining map)}\label{app:catalog}
% ========================================================================
This appendix collects the ``physics library of candidate models'' from the main text into one table:
for each physics family, its physical mechanism, the AI capability it might correspond to, and how far
AI has already used it. Maturity has three levels: \emph{mature} (representative work exists),
\emph{in progress} (some exploration, not yet mainstream), and \emph{potential} (a nearly blank,
high-potential area). This paper mines the ``integrable systems'' entry.

\begin{table}[h]
\centering
\footnotesize
\setlength{\tabcolsep}{4pt}
\renewcommand{\arraystretch}{1.25}
\begin{tabularx}{\textwidth}{@{}>{\raggedright\arraybackslash}p{2.4cm}
  >{\raggedright\arraybackslash}X >{\raggedright\arraybackslash}X
  >{\centering\arraybackslash}p{1.6cm} >{\raggedright\arraybackslash}X@{}}
\toprule
Physics family & Mechanism & Corresponding AI capability & Maturity & Status / representative direction \\
\midrule
Ising / spin glass & energy landscape, relaxation to the ground state & associative memory, denoising, energy models, generation & mature & Hopfield, Boltzmann machines; 2024 Nobel Prize in Physics \\
Non-equilibrium thermo.\ / diffusion & noise then denoise, a diffusion process & generation (image, audio, \dots) & mature & diffusion / score-based models \\
Symmetry / group equivariance & invariance and equivariance & geometric data, molecules, graphs & mature & equivariant networks, geometric deep learning \\
Hamiltonian / oscillator & conserved state along time evolution & long-range sequence modeling & in progress & coRNN, UnICORNN, LinOSS (the SSM line) \\
Renormalization group / coarse-graining & integrating out detail scale by scale & multiscale hierarchy, long-term memory & in progress & RG and deep-learning theory; RGMem, etc. \\
Tensor networks / quantum many-body & matrix-product states compressing a wavefunction & compression, classification, generation & in progress & MPS/TT classification and generation (niche) \\
Chaos / reservoir computing & dynamical systems at the edge of chaos & sequence prediction (nearly training-free) & in progress & reservoir computing / ESN \\
\textbf{Integrable systems} & \textbf{infinitely many conserved quantities, solitons, exact reversibility} & \textbf{length extrapolation, conserved / reversible sequences, crosstalk-free channels} & \textbf{potential} & \textbf{the first entry mined here (integrable / reversible)} \\
Lindblad / open quantum systems & evolution driven by dissipation and noise together & to be explored (noisy reasoning?) & potential & blank \\
KPZ / surface growth & random growth with scaling laws & to be explored (generation? scaling structure?) & potential & blank \\
Hubbard / Heisenberg & strongly correlated many-body & to be explored & potential & blank \\
\bottomrule
\end{tabularx}
\caption{The ``mining map'' of physics candidate model families. The rows marked ``potential'' are the
targets this project cares about most; we claim the integrable-systems entry first.}
\label{tab:catalog}
\end{table}

% ========================================================================
\section{Soliton-amplitude content: extraction protocol and self-check}\label{app:metric}
% ========================================================================
The amplitude content $\amp(x)$ is computed by a fixed protocol, applied the \emph{same} way to the
ground truth and to model predictions:

\begin{enumerate}
\item \textbf{Evolve to separation, read run lengths.} Pad the state with enough zeros on the right and
      evolve it under the finite-carrier dynamics until the solitons fully separate; the criterion is
      that the multiset of run lengths stays \emph{unchanged} after $\Delta$ more steps. Then read the
      length of each run of consecutive $1$s to get the sorted multiset. The step count needed for
      separation starts at $2|x|+60$ and doubles until it stabilizes.
\item \textbf{Padding amount (safety bound).} Under an unbounded carrier the rightmost ball moves right
      by at most ``total ball count'' cells per step, so the padding is set to
      $(\text{steps})\times(\text{total balls})+|x|+10$, which keeps the state within the right boundary.
\item \textbf{Handling garbage states.} A model prediction can be an illegal state (wrong ball count, no
      soliton structure). To avoid the pathological cost of a state that ``never separates'', we return
      the current best estimate after at most $4$ doublings. This is applied to all entrants and
      introduces no bias.
\item \textbf{No bias toward the structural model (key).} The extractor is \emph{discriminating}, not
      uniformly lenient: it gives $0\%$ to garbage output (generic models), $\sim40\%$ to partially
      correct output (the fixed rule), and $100\%$ to fully correct output (the structural model). Three
      clear levels show that it measures ``whether the soliton structure was actually preserved'', not a
      mechanical favoring of one side.
\end{enumerate}

\paragraph{Self-check coverage.} We verified that $\amp$ is \emph{conserved} under the finite-carrier
dynamics and that the system is \emph{reversible} (mirror method), across capacities $K=2,4,6$ with
several random multi-soliton configurations each (including cases with amplitude above $K$, which trigger
pass-through); the textbook $\{5,3,1\}$ configuration keeps amplitude content $\{5,3,1\}$ and ball count
$9$ under the dynamics.
% Maturity: hard numerical (self-check passes).

% ========================================================================
\section{Negative results from scouting (an honest memo)}\label{app:negatives}
% ========================================================================
\textbf{TODO: three parts.}
(i) \emph{Plain BBS / multi-soliton leak}: without a capacity limit the structural model's emit gate is
trivial (the fixed \texttt{emit}$=1-$cell already scores $100\%$), so we moved to finite capacity to get
a nontrivial, genuinely learned gate;
(ii) \emph{Colored BBS rule is not reversible}: the naive colored carrier rule is only $\sim86\%$
reversible under the natural mirror inverse, so it is not a canonical integrable rule and needs a crystal
/ combinatorial-$R$ construction; we defer the colored case to a future Tier 2;
(iii) \emph{Two Toda attempts} (action-angle, Neural-Lax): on a continuous system the free-form model
already solves it and the structure has no gap.

% ========================================================================
% Bibliography
% ========================================================================
\bibliographystyle{unsrt}
\bibliography{references}

\end{document}
